\documentclass[10pt, oneside]{article}   	% use "amsart" instead of "article" for AMSLaTeX format
\usepackage{geometry}                		% See geometry.pdf to learn the layout options. There are lots.
\geometry{letterpaper}                   		% ... or a4paper or a5paper or ... 
%\geometry{landscape}                		% Activate for rotated page geometry
%\usepackage[parfill]{parskip}    		% Activate to begin paragraphs with an empty line rather than an indent
\usepackage{graphicx}				% Use pdf, png, jpg, or eps§ with pdflatex; use eps in DVI mode
								% TeX will automatically convert eps --> pdf in pdflatex		
\usepackage{amssymb}
\usepackage{mhchem}
\usepackage{hyperref}

\title{Module 13, Part 2}
\author{Mark Peever\\ \texttt{mpeever@gmail.com}}
\date{March 27, 2026}

\begin{document}
\maketitle

\section*{Objectives}
\marginpar{0 minutes}
Refer to \href{https://drive.google.com/file/d/1fjlyN9yBxzAyjC-e-0vPJIUG56TnYqO6/view?usp=sharing}{Module 13 Notes}.\\

By the end of this class, the students should be able to\ldots
\begin{itemize}
\item describe Enthalpy
\item describe Hess's Law
\item describe Entropy
\item describe Gibbs Free Energy
\end{itemize}

\section*{Welcome \& Devotion}
\marginpar{5 minutes}
\begin{itemize}
\item have one student read \href{https://www.biblegateway.com/passage/?search=psalm\%20107\&version=LSB}{Psalm 107:17--22}
\end{itemize}


\section*{Review Enthalpy}
\marginpar{15 minutes}

\begin{itemize}
\item \emph{enthalpy} $H$ is energy absorbed into a substance
\item remember: $\Delta H > 0$ means endothermic
\item remember: $\Delta H < 0$ means exothermic
\end{itemize}

\begin{itemize}
\item it's important to notice that bonded atoms have \emph{lower} energy than unbonded atoms
\item so we measure enthalpy as ``energy absorbed" \emph{i.e.} $\Delta H > 0$
\item we calculate entropy as $\Delta H = H_{breaking} - H_{forming}$
\end{itemize}

\begin{itemize}
\item there is a baseline enthalpy ``enthalpy of formation" ($\Delta H_{f}^{0}$)
\begin{enumerate}
\item $\Delta H_{f}^{0} = 0$ for any element (in its elemental state)
\item $\Delta H_{f}^{0} = 0$ for any homonuclear diatomic
\item $\Delta H_{f}^{0} $ is \emph{some number} for any compound (see Table 13.2, p. 431)
\end{enumerate}
\item we can calculate enthalpy as $\Delta H^{0} = \Sigma \Delta H_{f} ^{0} (products) - \Sigma \Delta H_{f} ^{0} (reactants)$ 
\end{itemize}

\section*{Hess's Law}
\marginpar{10 minutes}

\begin{itemize}
\item Hess's Law says Enthalpy (and Entropy) depends on state, not on the path to get to it
\item so we don't need to know every step in a reaction to be able to calculate $\Delta H$ for a reaction or chain of reactions
\end{itemize}

\section*{Entropy}
\marginpar{20 minutes}

\begin{itemize}
\item Entropy ($S$) is a measure of disorder in a system
\item entropy applies to every physical system, but we measure it differently in different systems
\begin{enumerate}
\item homogeneous systems have higher entropy than heterogenous systems
\item gasses have higher entropy than liquids, liquids have higher entropy than solids
\item we can estimate entropy with numbers of particles (more particles means more chaos)
\end{enumerate}
\item \emph{e.g.} which side has higher entropy? \ce{CH4 (g) + 2O2 (g) -> 2H2O (g) + CO2 (g) + \Delta H}
%\item \emph{e.g.} which side has higher entropy? \ce{H2SO4 (aq) + 2H2O (l) -> 2H3O+ (aq) + SO4^{2-} (aq)}
\item \emph{e.g.} which side has higher entropy? \ce{H2SO4 (aq) + NaOH (aq) -> 2H2O (l) + Na2SO4 (aq)}
\item the Second Law says: $\Delta S_{universe} \ge 0$
\item discuss the Second Law at length and relate to various things
\end{itemize}

\section*{Gibbs Free Energy}
\marginpar{20 minutes}

\begin{itemize}
\item reactions want to be exothermic
\item reactions want to increase entropy
\item $\Delta G = \Delta H - T \cdot \Delta S$
\item $\Delta G$ allows us to predict the spontaneity of a rection: if $\Delta G < 0$, a reaction is spontaneous
\item notice we can change $\Delta G$ by doing things like changing temperature!
\end{itemize}


\section*{Questions for me}
\marginpar{5 minutes}

\section*{Assignment}
\marginpar{5 minutes}
\begin{itemize}
\item Module 12 Practice Problems: p. 416 \#1--10 due 2026-04-10
\item Module 13 Review Problems: p. 456 \# 1--10 (not to be turned in)
\item Module 13 Practice Problems: pp. 457--458 \#1--10 due 2026-04-10
\end{itemize}



\end{document}  